\chapter{Introducción} \label{sec:Introduccion}

\section{Contexto y motivación}

La adopción de Kubernetes como plataforma de orquestación de contenedores se ha consolidado en la industria tecnológica, pero su viabilidad operativa depende en gran medida del componente de red que conecta los servicios dentro del clúster. Kubernetes no incluye por defecto un mecanismo de conectividad entre contenedores: delega esa función en un complemento intercambiable denominado Container Network Interface (CNI), cuya selección recae en el equipo técnico de cada organización.

Los equipos deben escoger entre múltiples complementos ---Calico, Cilium, Flannel y Antrea, entre otros--- sin contar con criterios de comparación objetivos ni datos de rendimiento medidos en condiciones representativas. Esta carencia conduce a tres problemas interrelacionados: configuraciones de red inseguras por la elección de complementos que no aplican políticas de control de acceso en su plano de datos, sobredimensionamiento de recursos computacionales por desconocimiento del costo real de cada complemento en CPU y memoria, y una exposición silenciosa a vulnerabilidades que el equipo operativo no detecta hasta que afectan a los usuarios.

La pregunta de investigación que motiva este trabajo es la siguiente: ¿cómo pueden las organizaciones comparar complementos de red y automatizar políticas de comunicación segura para optimizar recursos sin degradar el rendimiento?

De esta pregunta principal se derivan cuatro sub-preguntas específicas que estructuran la evaluación empírica:

\begin{enumerate}
    \item ¿Cuáles son las diferencias en velocidad de respuesta ---throughput, latencia y retransmisiones TCP--- entre los plugins CNI evaluados bajo condiciones controladas de carga?
    \item ¿Cómo impacta cada CNI en el consumo de recursos del nodo ---CPU y memoria RAM--- durante sesiones de red de alta intensidad?
    \item ¿Qué tan efectivo es cada CNI para aplicar Network Policies de seguridad, y cuál es el overhead que introduce sobre el rendimiento?
    \item ¿En qué medida la diferencia de eficiencia entre plugins permite reducir el sobredimensionamiento de infraestructura a escala de clúster?
\end{enumerate}

La evaluación empírica de complementos de red permite fundamentar esas decisiones con datos medidos, no con suposiciones. Este trabajo proporciona esa evidencia para cuatro CNIs representativos, junto con una metodología replicable por cualquier equipo que necesite realizar la misma comparación. Para operacionalizar esta evaluación, el proyecto provee dos entregables técnicos principales: un banco de pruebas de red (testbed) reproducible mediante infraestructura como código (IaC) para la obtención automática de métricas, y una herramienta interactiva recomendadora en formato de aplicación web de página única (SPA). Estos componentes traducen los datos empíricos recopilados del plano de datos en decisiones estructuradas de configuración y dimensionamiento para las organizaciones.

\section{Planteamiento del problema} \label{sec:Planteamiento}

Kubernetes se ha convertido en el estándar de facto para desplegar aplicaciones modernas en contenedores, pero el rendimiento, la seguridad y la eficiencia de un clúster no dependen solo del orquestador. Dependen del complemento de red que conecta los servicios entre sí. Esta responsabilidad recae en los plugins CNI, que determinan cómo se enrutan los paquetes, qué tan rápido viajan y qué comunicaciones están permitidas o bloqueadas.

El problema es que no existe un criterio objetivo y ampliamente adoptado para elegir entre los distintos CNIs disponibles. Las decisiones de selección suelen apoyarse en documentación dispersa, recomendaciones informales o en el plugin configurado por defecto ---en el caso de K3s, Flannel---, sin datos de rendimiento medidos en condiciones representativas. El resultado concreto es triple: configuraciones de red inseguras, porque algunos complementos no aplican en el tráfico real las políticas que el operador cree activas; clústeres sobredimensionados, porque sin conocer el costo real de CPU y RAM de cada CNI se agregan recursos a modo de margen de seguridad; y fallos que solo se detectan cuando ya afectan al usuario final, porque la red interna del clúster carece de observabilidad.

El problema se agrava porque Kubernetes provee los mecanismos necesarios ---NetworkPolicies, autoescalado, observabilidad---, pero su aplicación efectiva requiere conocimientos especializados y pruebas sistemáticas que la mayoría de equipos no tiene capacidad de realizar. En entornos donde la disponibilidad es crítica, operar sobre supuestos en lugar de evidencias convierte la red en el punto débil del sistema.

Esta investigación parte de la necesidad de comparar CNIs con datos objetivos y de automatizar la validación de políticas de red, para que las organizaciones puedan garantizar seguridad, rendimiento y eficiencia sin incrementar la complejidad operativa. Para resolver esta brecha, el proyecto propone dos entregables técnicos: un banco de pruebas de red (testbed) reproducible basado en infraestructura como código (IaC) para la recolección automática de telemetría y una aplicación web interactiva (SPA) que implementa el algoritmo MCDA. Estos instrumentos reducen la complejidad operativa al traducir las mediciones del plano de datos en decisiones de configuración directa y dimensionamiento correcto de recursos.

\section{Justificación} \label{sec:Justificacion}

La adopción de Kubernetes como estándar de facto para el despliegue de aplicaciones en contenedores ha trasladado una decisión crítica al equipo técnico de cada organización: la selección del complemento de red (CNI) que conecta los servicios dentro del clúster. Esta decisión es de naturaleza telemática ---determina cómo se enrutan los paquetes, con qué latencia y throughput viajan, qué comunicaciones se permiten y cuántos recursos consume el plano de red---, pero en la práctica se toma sin evidencia medida, apoyándose en documentación dispersa o en el complemento configurado por defecto~\cite{ref29}. Este trabajo se justifica precisamente por esa brecha: convierte una elección basada en intuición en una decisión fundamentada en datos reproducibles.

\subsection{Pertinencia telemática}

El aporte se ubica en el núcleo de la ingeniería telemática: la convergencia entre redes de datos y servicios distribuidos. Tres dimensiones telemáticas estructuran la solución. La primera es la \textbf{Calidad de Servicio (QoS)}: la medición controlada de throughput, latencia de establecimiento y retransmisiones TCP entre nodos cuantifica el comportamiento real de la red del clúster bajo carga. La segunda es la \textbf{observabilidad de red}: la instrumentación con Prometheus y Grafana sobre el agente CNI hace visible el tráfico este-oeste interno, un plano que las herramientas tradicionales de monitoreo no cubren. La tercera es la \textbf{micro-segmentación bajo Zero Trust}: la automatización y validación de NetworkPolicies traduce el principio de mínimo privilegio en reglas verificables sobre el plano de datos. En conjunto, el proyecto no evalúa un producto de software genérico, sino el comportamiento telemático de la infraestructura de comunicaciones de un clúster.

\subsection{Impacto en las organizaciones}

El beneficio para las organizaciones que sigan esta guía de implementación es directo y medible. Conocer el costo real en CPU y memoria de cada CNI por unidad de tráfico permite dimensionar la infraestructura con base en evidencia y reducir el \emph{sobredimensionamiento preventivo} ---la práctica de adquirir capacidad ``por si acaso'' que genera gasto innecesario en infraestructura y energía~\cite{ref7,ref8}. La validación automatizada de políticas de red cierra brechas de seguridad que suelen permanecer ocultas hasta que un incidente las revela, sin elevar la complejidad operativa. Y la entrega de un método replicable ---infraestructura como código, \emph{benchmarks} automatizados y un modelo de decisión---  permite que equipos sin especialización profunda en redes de Kubernetes realicen la misma comparación en su propio entorno. El impacto se resume en tres ejes: decisiones de arquitectura basadas en datos, reducción de costos operativos y una postura de seguridad verificable.

\subsection{Aporte académico y entregables}

Frente a los estudios comparativos existentes, que suelen reportar mediciones aisladas de rendimiento, este trabajo integra la evaluación de QoS, la seguridad y la eficiencia de recursos en un único modelo de decisión multicriterio parametrizable por perfil de uso. El resultado tangible no es solo el documento, sino una \textbf{herramienta}: un testbed reproducible y un prototipo recomendador que expone, para cada perfil organizacional, qué CNI conviene y por qué. Esta orientación hacia un producto verificable ---y no únicamente hacia un informe--- es coherente con la modalidad de monografía orientada a la solución de un problema tecnológico mediante modelos, prototipos y productos~\cite{ref29}.

\subsection{Viabilidad}

El proyecto es viable técnica, económica y legalmente. La totalidad del \emph{stack} ---Kubernetes (K3s), Calico, Cilium, Flannel, Antrea, Prometheus, Grafana, iperf3--- se distribuye bajo licencias permisivas (Apache~2.0, MIT, BSD o GPL) que habilitan su uso académico sin restricciones. La infraestructura se aprovisiona bajo demanda en la nube pública con un costo acotado y controlado mediante automatización, y no se manipulan datos personales ni información sensible, por lo que no aplican consideraciones de protección de datos.

\section{Objetivos} \label{sec:Objetivos}

\subsection{Objetivo General}

Desarrollar un modelo de evaluación, selección y validación de plugins CNI en Kubernetes, compuesto por una guía metodológica, un prototipo de pruebas de laboratorio y una lista de verificación.

\subsection{Objetivos Específicos}

\begin{enumerate}
  \item Evaluar cuantitativamente el rendimiento de los CNI disponibles para Kubernetes (Calico, Cilium, Flannel y Antrea) mediante pruebas controladas de latencia, velocidad de transferencia y uso de recursos, para identificar sus fortalezas y limitaciones en escenarios reales.
  \item Evaluar cómo la configuración segura y automatizada de las políticas de red adaptadas a cada plugin CNI reduce la superficie de ataque y la probabilidad de incidentes de seguridad en clústeres Kubernetes, con el fin de determinar su capacidad de reducir vulnerabilidades y prevenir incidentes de seguridad en entornos de laboratorio representativos de producción.
  \item Desarrollar, como parte del modelo propuesto, un conjunto de criterios y métricas que permitan la comparación objetiva entre diferentes CNI para Kubernetes mediante la definición de umbrales de rendimiento y factores de ponderación, para facilitar la toma de decisiones técnicas basadas en evidencia cuantificable.
  \item Implementar un prototipo funcional que recomiende y configure el CNI más adecuado según las necesidades específicas de cada organización, para simplificar el proceso de implementación y apoyar la definición de configuraciones recomendadas en entornos de prueba.
\end{enumerate}

\section{Alcance del trabajo}

Esta tesis tiene como objetivo desarrollar un modelo de evaluación, selección y validación de complementos CNI en clústeres Kubernetes. El trabajo evalúa Flannel, Calico, Cilium y Antrea en un clúster K3s hospedado en la nube de DigitalOcean, midiendo la velocidad de transferencia, la latencia, el uso de CPU, de memoria y la efectividad de las políticas de seguridad. No obstante, esta investigación solo considera configuraciones de red superpuestas (overlay). El estudio excluye despliegues en servidores físicos (bare-metal) o enrutamiento nativo debido a las restricciones de la red virtual privada.

\section{Estructura del documento}

Este capítulo introductorio ha presentado el contexto y la motivación, el planteamiento del problema, la justificación, los objetivos y el alcance del trabajo. El \cref{cap:marco_teorico} desarrolla el marco teórico-conceptual y el \cref{sec:EstadoArte} revisa la literatura sobre la evaluación de CNIs, la seguridad Zero Trust y los modelos de decisión. Los \cref{cap:metodologia,sec:Metodos,cap:herramienta} detallan la metodología de trabajo, el diseño de la solución y su implementación y despliegue. Finalmente, los \cref{sec:Resultados,sec:Discusion,sec:Conclusiones,sec:Recomendaciones,sec:TrabajosFuturos} documentan los resultados experimentales, la discusión de implicaciones, las conclusiones, las recomendaciones y el trabajo futuro.
